% ============================================================================
% Fichier  : partie5/P05_C18_confinement_retablissement.tex
% Titre    : Confinement, Éradication, Rétablissement
% Auteur   : MINKA MI NGUIDJOÏ Thierry Emmanuel
% ============================================================================
% Sources : ISO/IEC 27035 ; NIST SP 800-61r2. Opération MAYO.
% ============================================================================

\chapter{Confinement, Éradication, Rétablissement}
\label{chap:confinement}
\chaptermark{Confiner, éradiquer, rétablir}

\epigraph{%
    « Contenir sans détruire la preuve, éradiquer sans oublier une porte,\\
    rétablir sans ramener l'attaquant. »%
}{Principe du cours}

\niveauChapitre{Intermédiaire à expert}{Chapitre~\ref{chap:cycle}.}

\begin{boxobjectifs}
\begin{itemize}
    \item Choisir une stratégie de confinement adaptée.
    \item Éradiquer complètement la présence de l'attaquant.
    \item Rétablir à partir de sauvegardes sûres.
    \item Vérifier qu'il n'y a pas de rechute.
    \item Décider, ou non, de payer une rançon.
\end{itemize}
\end{boxobjectifs}

\bigskip

% ============================================================================
\section*{Question directrice : comment remettre LiZbiZ en marche sans rouvrir la porte~?}
% ============================================================================

Après la qualification et la capture des preuves, il faut agir~: contenir,
éradiquer, rétablir. Chaque étape a ses pièges. Ce chapitre les parcourt, et
aborde la question difficile du paiement de la rançon.

\medskip
\textbf{Plan du chapitre.}
\begin{enumerate}
    \item Le confinement (section~\ref{sec:c18-confiner}).
    \item L'éradication (section~\ref{sec:c18-eradiquer}).
    \item Le rétablissement (section~\ref{sec:c18-retablir}).
    \item Faut-il payer~? (section~\ref{sec:c18-payer}).
\end{enumerate}

\separateur

% ============================================================================
\section{Le Confinement}
\label{sec:c18-confiner}
% ============================================================================

\begin{tableaucours}{L{3.6cm}L{4.6cm}L{4.0cm}}{%
    \tableauHeader{Stratégie} &
    \tableauHeader{Principe} &
    \tableauHeader{Quand}
}
Isolement réseau & Couper la machine du réseau sans l'éteindre & Préserve la preuve \\
Segmentation & Fermer un segment & Limiter la propagation \\
Blocage de comptes & Désactiver les comptes compromis & Couper l'accès \\
Arrêt (dernier recours) & Éteindre & Si la propagation l'exige \\
\end{tableaucours}
\vspace{-0.8em}
\begin{center}{\small\itshape Tableau~18.1. Stratégies de confinement}\end{center}

On privilégie l'isolement réseau, qui préserve la mémoire et la preuve
(Chapitre~\ref{chap:memoire}).

% ============================================================================
\section{L'Éradication}
\label{sec:c18-eradiquer}
% ============================================================================

\begin{boxprocedure}
\textbf{Éradiquer}
\begin{enumerate}
    \item Supprimer les accès de l'attaquant (comptes créés, clés déposées).
    \item Supprimer les persistances (services, tâches, clés de démarrage).
    \item Corriger la faille d'entrée (sinon l'attaquant revient).
    \item Changer les secrets exposés (mots de passe, clés).
    \item Vérifier l'éradication avant de rétablir.
\end{enumerate}
\end{boxprocedure}

\begin{boxalert}
\textbf{Erreur fréquente.} Supprimer le rançongiciel visible et croire la
machine propre.

\smallskip
\textbf{Règle~:} l'attaquant a souvent plusieurs portes (comptes, services,
accès distant). On éradique à partir de la chronologie forensique, qui liste
\emph{toutes} ses traces, pas seulement la charge finale.
\end{boxalert}

% ============================================================================
\section{Le Rétablissement}
\label{sec:c18-retablir}
% ============================================================================

On rétablit à partir de sauvegardes \termecle{vérifiées} et \termecle{hors
ligne} (une sauvegarde en ligne a pu être chiffrée elle aussi). On restaure
par étapes, en surveillant étroitement, prêt à reconfiner en cas de rechute.
La supervision renforcée (Partie~II) accompagne le retour à la normale.

% ============================================================================
\section{Faut-il Payer~?}
\label{sec:c18-payer}
% ============================================================================

\begin{boxscenario}{Opération MAYO. LiZbiZ ne paie pas}
\textbf{Constatation.} La rançon demandée s'élève à 25~millions de FCFA.
LiZbiZ dispose de sauvegardes hors ligne du serveur de suivi, vieilles de
cinq jours. Les données exfiltrées, elles, sont déjà perdues.

\medskip
\textbf{Analyse} (confiance élevée). Payer ne garantit ni la restitution ni
la non-divulgation, finance le crime, et peut être juridiquement risqué.
LiZbiZ choisit de ne pas payer~: elle restaure depuis les sauvegardes
(perte de cinq jours de données de suivi, reconstituées), notifie l'ANTIC et
les clients concernés par l'exfiltration. La décision est documentée.
\end{boxscenario}

\begin{boxCRO}
\textbf{La décision de payer, sous l'angle CRO.} Payer n'améliore pas la
\cro{-}{-}{-}~: les données exfiltrées restent exposées (C), rien ne garantit
l'intégrité des données « rendues » (F), et le paiement affaiblit la position
devant l'autorité (O). La préparation (sauvegardes) est la vraie réponse.
\end{boxCRO}

\begin{boxmetier}{Responsable de la réponse à incident et RSSI}
\textbf{La meilleure réponse au rançongiciel est antérieure à l'attaque.}
Des sauvegardes hors ligne testées rendent la question du paiement presque
théorique. Le RRI prépare cette réponse à froid.
\end{boxmetier}

\separateur

\begin{boxresume}
\textbf{Ce qu'il faut retenir de ce chapitre~:}
\begin{itemize}
    \item \textbf{Confiner} en isolant, pour préserver la preuve.
    \item \textbf{Éradiquer} toutes les portes, pas seulement la charge
          visible.
    \item \textbf{Corriger la faille d'entrée} avant de rétablir.
    \item \textbf{Rétablir} depuis des sauvegardes vérifiées hors ligne.
    \item \textbf{Payer} n'est pas une solution~; la préparation l'est.
\end{itemize}
\end{boxresume}

\begin{boxinfo}
\textbf{Points de contrôle de l'élève.} Je sais~: choisir une stratégie de
confinement~; éradiquer complètement~; rétablir sûrement~; vérifier l'absence
de rechute~; argumenter sur le paiement.
\end{boxinfo}

\bigskip

\begin{boxexo}[title={\ding{45}\quad Exercice~18.1. Confiner \niveaubase}]
Pourquoi préférer l'isolement réseau à l'extinction~?
\end{boxexo}

\begin{boxexo}[title={\ding{45}\quad Exercice~18.2. Portes \niveaubase}]
Citez trois « portes » qu'un attaquant peut laisser.
\end{boxexo}

\begin{boxexo}[title={\ding{45}\quad Exercice~18.3. Sauvegardes \niveaubase}]
Pourquoi rétablir depuis des sauvegardes hors ligne~?
\end{boxexo}

\begin{boxexo}[title={\ding{45}\quad Exercice~18.4. Éradication incomplète \niveaumoyen}]
Pourquoi supprimer seulement le rançongiciel visible ne suffit-il pas~?
\end{boxexo}

\begin{boxexo}[title={\ding{45}\quad Exercice~18.5. Faille d'entrée \niveaumoyen}]
Que se passe-t-il si l'on rétablit sans corriger la faille d'entrée~?
\end{boxexo}

\begin{boxexo}[title={\ding{45}\quad Exercice~18.6. Payer \niveaumoyen}]
Donnez trois raisons de ne pas payer une rançon.
\end{boxexo}

\begin{boxexo}[title={\ding{45}\quad Exercice~18.7. CRO \niveaumoyen}]
Analysez la décision de payer sous l'angle CRO.
\end{boxexo}

\begin{boxexo}[title={\ding{45}\quad Exercice~18.8. MAYO \niveaumoyen}]
Reconstituez la décision de LiZbiZ~: pourquoi ne pas payer, et comment
rétablir~?
\end{boxexo}

\begin{boxexo}[title={\ding{45}\quad Exercice~18.9. Plan de rétablissement \niveauexpert}]
Rédigez le plan de rétablissement de LiZbiZ après l'éradication.
\end{boxexo}

\begin{boxexo}[title={\ding{45}\quad Exercice~18.10. Politique rançongiciel \niveauexpert}]
Rédigez la politique de LiZbiZ sur le paiement des rançons, fondée sur la
préparation.
\end{boxexo}

\bigskip

\begin{boxinfo}
\textbf{Pour aller plus loin.}
\begin{itemize}
    \item ISO/IEC 27035~\cite{iso27035}~; NIST SP~800-61r2~\cite{nist80061r2}.
\end{itemize}
\end{boxinfo}

\begin{boxinfo}
\textbf{Transition.} L'incident traité, il reste à en tirer les leçons~: la
Partie~VI s'ouvre sur le retour d'expérience
(Chapitre~\ref{chap:retex}).
\end{boxinfo}
